\documentclass[11pt]{article}
\usepackage[T1]{fontenc}
\usepackage[utf8]{inputenc}
\usepackage{lmodern,amsmath,amssymb,amsthm,booktabs,microtype}
\usepackage[margin=28mm]{geometry}
\usepackage[colorlinks=true,linkcolor=blue,citecolor=blue,urlcolor=blue]{hyperref}
\newtheorem{theorem}{Theorem}[section]
\newtheorem{lemma}[theorem]{Lemma}
\newtheorem{proposition}[theorem]{Proposition}
\newcommand{\R}{\mathbb R}
\newcommand{\g}{\mathfrak g}
\newcommand{\Lie}{\operatorname{Lie}}
\newcommand{\ad}{\operatorname{ad}}
\title{Universal identities for Euclidean mechanical Lie brackets}
\author{Kaiwen Jin\\\small\href{mailto:yui.ui.siki@gmail.com}{yui.ui.siki@gmail.com}}
\date{October 11, 2026}

\begin{document}
\maketitle
\begin{abstract}
We prove that Euclidean kinetic energy and arbitrary smooth potentials,
considered in all dimensions, satisfy no universal Lie identities beyond
those of the free polynomially graded Lie algebra of classical mechanics.
Polynomial potentials already suffice to separate every nonzero element.
The proof reduces positive momentum degree to degree zero: a sufficiently
large binary-tree potential distinguishes words in degree one, while in
higher degree a free Lie projection and a retraction recover a nonzero
multiple of the original element. Symmetrization in a free pre-Lie algebra
and finite Taylor jets then separate the potential part. This answers the
Euclidean mechanical conjecture of McLachlan and Murua. Exact computations
through order eighteen provide supporting finite verification.
\end{abstract}

\section{Introduction}

For a mechanical Hamiltonian
\[
 H(q,p)=T(p)+V(q),\qquad T(p)=\tfrac12\lvert p\rvert^2,
\]
the Poisson brackets of kinetic and potential energy determine the
order conditions of splitting integrators. Quadratic dependence on momentum
forces identities absent from an unrestricted free Lie algebra. For example,
three successive brackets with a potential annihilate the kinetic energy.
The question is whether a constant Euclidean kinetic metric imposes further
identities that hold for every potential and every dimension.

McLachlan and Murua \cite{MM} formulated this question using the free
polynomially graded Lie algebra of classical mechanics. Their structural
theorem identifies its potential part with a free commutative
nonassociative algebra, and its positive-degree part with a free Lie algebra
on explicit generators. They also established faithfulness on the potential
part by relating it to rooted trees and pre-Lie symmetrization. The
corresponding assertion for all momentum degrees remained their
Conjecture~1. The survey of Blanes, Casas and Murua \cite[\S3.1]{BCM} retains
the conjectural equality of the Euclidean and general mechanical
order-condition reductions and reports prior checks through order twenty.
Finite-order rank calculations alone cannot establish this equality in all
orders.

Here is the precise free object and the statement we prove. Work over $\R$.
Let $\mathfrak P$ be the class of Lie algebras $L=\bigoplus_{m\ge0}L_m$ with
\begin{equation}\label{eq:grading}
 [L_n,L_m]\subseteq L_{n+m-1}\quad(n+m>0),
 \qquad [L_0,L_0]=0.
\end{equation}
Let $F$ be free in this class on $A$ of degree two and $B$ of degree zero;
every degree-preserving assignment of these generators extends uniquely to
a homomorphism. For $d\ge1$ and a smooth real potential $V$ on $\R^d$, write
\[
 \Phi_{d,V}:F\longrightarrow C^\infty_{\mathrm{pol}}(\R^{2d}),
 \qquad A\longmapsto\tfrac12\lvert p\rvert^2,
 \qquad B\longmapsto V(q).
\]
The target consists of smooth functions polynomial in $p$, graded by
homogeneous momentum degree, with the convention
\begin{equation}\label{eq:poisson}
 \{f,h\}=\sum_{i=1}^d(f_{q_i}h_{p_i}-f_{p_i}h_{q_i}).
\end{equation}

\begin{theorem}\label{thm:main}
The Euclidean mechanical realizations have zero joint kernel:
\begin{equation}\label{eq:main}
 \bigcap_{d\ge1}\ \bigcap_{V\in C^\infty(\R^d;\R)}
       \ker\Phi_{d,V}=\{0\}.
\end{equation}
The intersection is already zero when the potentials are restricted to real
polynomials.
\end{theorem}

Thus every nonzero finite Lie expression is detected by some dimension and
some polynomial potential. The quantifiers in \eqref{eq:main} are those of
the Euclidean conjecture and AIM problem~519; injectivity for an individual
fixed potential is not asserted.

The new step is an algebraic degree-lowering argument. A putative universal
identity belongs to a graded ideal. Bracketing with a suitable modified
potential lowers its momentum degree without making it zero. The probe may
have a much larger order than the expression being tested, which is permitted
by the universal question. In degree one a large binary tree makes the
ordered sequence of multiplications recoverable. In higher degree the probe
acts, after discarding generators with larger tree labels, like replacing one
independent kinetic letter by a fresh letter. Sending that letter back
recovers the original expression multiplied by its kinetic-letter count.
Every nonzero graded ideal must consequently meet the potential part, where
faithfulness closes the proof.

Section~\ref{sec:lowering} proves this reduction from the established free
decomposition. Section~\ref{sec:potentials} gives explicit tree and Taylor-jet
arguments for the potential part and completes the theorem.
Appendix~\ref{sec:computations} describes the finite computations. All Lie
elements and tree combinations in the paper are finite sums; no completed
Lie algebra is used.

\section{Lowering momentum degree}\label{sec:lowering}

Put $X=F_0$ and define
\begin{equation}\label{eq:star}
 x*s=[s,[x,A]]\qquad(x,s\in X).
\end{equation}
Jacobi and $[X,X]=0$ show that this product is commutative. Theorem~1 of
McLachlan and Murua \cite{MM}, together with its Hall-basis proof in
\S2.3, gives the vector-space decomposition
\begin{equation}\label{eq:decomposition}
 F=X\oplus\g,\qquad
 \g=\bigoplus_{m\ge1}F_m
    =\Lie_{\mathrm{free}}(A,\{y_s:s\in\mathcal B\}),
 \qquad y_s=[s,A].
\end{equation}
Here $\mathcal B$ is the basis of all unordered full binary products in $B$.
With \eqref{eq:star}, $X$ is the free \emph{nonunital} commutative
nonassociative algebra on $B$. The specified generators in
\eqref{eq:decomposition}, rather than abstract freeness alone, will be used
below. If $|s|$ denotes the number of leaves, then
\[
 |s|\ge1,\qquad |s*t|=|s|+|t|,
 \qquad [x,A]=y_x,\quad [x,y_s]=x*s.
\]
For general $x\in X$, the notation $y_x=[x,A]$ is extended linearly.

Count the occurrences of $A$ as an independent generator of $\g$;
the letters inside a label $y_s$ are not counted. A Lie monomial with $k$
such occurrences has momentum degree $k+1$, by
\eqref{eq:grading}. Freeness makes this a second grading, and hence
\begin{equation}\label{eq:Acount}
 F_1=\Lie_{\mathrm{free}}(\{y_s\}),\qquad
 F_m=\{\text{elements of $\g$ homogeneous of independent $A$ count $m-1$}\}.
\end{equation}
We use the injective associative realization of a free Lie algebra as
commutator polynomials in its generators. An algebraic justification is given
in Appendix~\ref{sec:associative}.

\begin{lemma}\label{lem:degreeone}
For every $0\ne u\in F_1$, there is $x\in\mathcal B$ such that
$[x,u]\ne0$.
\end{lemma}
\begin{proof}
Expand $u$ as a nonzero finite associative polynomial
$\sum_w c_w w$ in words $w=y_{s_1}\cdots y_{s_r}$.
Choose
\[
 M>\max_{c_w\ne0}\sum_{i=1}^r|s_i|,
 \qquad x\in\mathcal B,\quad |x|=M.
\]
Such a tree exists for every positive $M$, for instance a comb.
On $X$, $\ad_{y_s}=-L_s$, where $L_s(z)=s*z$. Thus the action of a word is
\begin{equation}\label{eq:wordaction}
 w(x)=(-1)^r s_1*(s_2*(\cdots*(s_r*x)\cdots)).
\end{equation}
These output trees distinguish the words. At the root, the branch containing
$x$ has at least $M$ leaves and the other branch has fewer than $M$.
Removing the latter recovers $s_1$. Repeat to recover $s_2,\ldots,s_r$, stopping
when the remaining tree has exactly $M$ leaves. Every proper subtree of $x$
has fewer leaves, so the stopping rule is unambiguous. Distinct associative
words therefore have distinct outputs in the basis $\mathcal B$.
Their coefficients, including the signs in \eqref{eq:wordaction}, cannot
cancel. Hence $[u,x]\ne0$, and antisymmetry proves the lemma.
\end{proof}

For higher momentum degree we must control the additional potential terms
created by this action. Fix $M\ge1$ and set
\[
 X_{>M}=\operatorname{span}\{s\in\mathcal B:|s|>M\},
 \qquad
 J_{>M}=\text{the Lie ideal of $\g$ generated by $y_s$ with $|s|>M$},
\]
and $N_{>M}=X_{>M}\oplus J_{>M}\subset F$.

\begin{lemma}\label{lem:stability}
The subspace $N_{>M}$ satisfies
\[
 [X_{>M},\g]\subseteq N_{>M},\qquad
 [N_{>M},\g]\subseteq N_{>M}.
\]
\end{lemma}
\begin{proof}
Induct on the length of a bracket expression $w$ in the free generators of
$\g$, uniformly for all $z\in X_{>M}$. For a generator,
$[z,A]=y_z\in J_{>M}$ and $[z,y_s]=z*s\in X_{>M}$, since $|s|\ge1$.
If $w=[v_1,v_2]$, Jacobi gives
\[
 [z,w]=[[z,v_1],v_2]+[v_1,[z,v_2]].
\]
The induction hypothesis splits each $[z,v_i]$ into an $X_{>M}$ term and a
$J_{>M}$ term. The former can be bracketed with the shorter expression $v_j$
using the same hypothesis; the latter remains in $J_{>M}$ by its ideal
property inside $\g$. This proves the first inclusion. Adding
$[J_{>M},\g]\subseteq J_{>M}$ gives the second.
\end{proof}

\begin{lemma}\label{lem:higher}
For every $m\ge2$ and $0\ne u\in F_m$, there is $x\in\mathcal B$ with
$0\ne[x,u]\in F_{m-1}$.
\end{lemma}
\begin{proof}
Choose $M$ strictly larger than every tree label in a finite free Lie
expression for $u$, and choose $|x|=M$. If there are no such labels take any
$M\ge1$. In particular $y_x$ is absent from $u$.
Define a free associative derivation, restricted to $\g$, by
\[
 \delta_x(A)=y_x,\qquad \delta_x(y_s)=0\quad(s\in\mathcal B).
\]
On a word it replaces one occurrence of $A$ by $y_x$ and sums over all such
positions. The commutator derivation rule shows that it preserves the free
Lie subalgebra.

Regard both $\delta_x$ and $\ad_x$ as derivations $\g\to F$. Their difference
is zero on $A$ and equals $x*s\in X_{>M}$ on $y_s$. Expanding a bracket by the
derivation rule and applying Lemma~\ref{lem:stability} inductively yields
\[
 [x,u]-\delta_xu\in N_{>M}.
\]
Both terms have momentum degree $m-1\ge1$, so the difference actually lies
in $N_{>M}\cap\g=J_{>M}$.

Let $\pi_M$ be the free Lie homomorphism that kills all $y_s$ with $|s|>M$
and retains $A$ and the other $y_s$. Its target is
$\g_{\le M}=\Lie_{\mathrm{free}}(A,\{y_s:|s|\le M\})$, identified with the
subalgebra on these generators. Its kernel is exactly $J_{>M}$: the quotient
$\g/J_{>M}$ and $\g_{\le M}$ are isomorphic by inverse
generator-preserving maps, by their universal properties. Consequently
\begin{equation}\label{eq:projection}
 \pi_M([x,u])=\delta_xu.
\end{equation}
The derivative uses only the old labels smaller than $M$ and the new label
of size $M$, so its image is retained.

Define $r_x:\g_{\le M}\to\g$ by sending $y_x$ to $A$ and fixing every other
generator, including $A$. For a bracket monomial without $y_x$,
$r_x\delta_x$ multiplies it by its independent $A$ count: this holds on each
generator and the derivation rule adds the counts under a bracket.
Equation~\eqref{eq:Acount} therefore gives
\begin{equation}\label{eq:euler}
 r_x\pi_M([x,u])=(m-1)u\ne0.
\end{equation}
Over $\R$ the scalar $m-1$ is nonzero, proving the claim. These auxiliary
maps are maps of ordinary free Lie algebras; no extension of $r_x$ to $F$
or grading preservation by $r_x$ is required.
\end{proof}

\begin{proposition}\label{prop:ideal}
Every nonzero graded Lie ideal of $F$ has a nonzero intersection with $X$.
\end{proposition}
\begin{proof}
If a nonzero graded ideal misses $X$, take its least occupied degree
$m\ge1$ and a nonzero homogeneous element $u$ in that degree.
Lemma~\ref{lem:degreeone} or Lemma~\ref{lem:higher} supplies a nonzero
$[x,u]$ of degree $m-1$ in the same ideal, a contradiction.
\end{proof}

\section{Separating potentials and completing the proof}\label{sec:potentials}

For completeness we make the degree-zero argument explicit, including the
passage from formal trees to actual potentials. Let $\mathcal U$ be the
vector space on finite nonempty unrooted simple trees. For two basis trees
$t,v$, define $t\diamond v$ by joining one vertex of $t$ to one vertex of
$v$ by a new edge and summing over all vertex pairs. Isomorphic outputs
are counted with multiplicity. This product is commutative. Freeness of $X$
defines a map
\[
 \psi:X\longrightarrow\mathcal U,\qquad
 \psi(B)=\text{one-vertex tree},\qquad
 \psi(x*s)=\psi(x)\diamond\psi(s).
\]

\begin{lemma}\label{lem:unrooted}
The map $\psi$ is injective.
\end{lemma}
\begin{proof}
Let $\mathcal R$ be the vector space on rooted trees, with
$a\triangleright b$ obtained by attaching the root of $b$ to each vertex of
$a$ and retaining the root of $a$. It is the free pre-Lie algebra on the
rooted singleton, by Chapoton and Livernet \cite[Theorem~1.9 and
Corollary~1.10]{CL}. Define $\rho:\mathcal U\to\mathcal R$ by summing over
all choices of a root. Then
\begin{equation}\label{eq:rooting}
 \rho(t\diamond v)=\rho(t)\triangleright\rho(v)
                         +\rho(v)\triangleright\rho(t).
\end{equation}
The first summand counts joined trees whose root lies in $t$, and the second
those whose root lies in $v$; the choices of root and attachment vertex agree
on both sides, including multiplicities from automorphisms.

Bergeron and Loday \cite[Theorem~3.1]{BL} prove that the operad morphism
$\mathrm{ComMag}\to\mathrm{preLie}$ induced by symmetrization is injective.
This gives the one-generator injection used here. Indeed, over $\R$,
coinvariants of a finite symmetric group are identified with invariants by
averaging; the inverse is the quotient map. Passing an equivariant injection
to coinvariants consequently remains injective. Tensoring each arity with
the one-dimensional generator space and taking the direct sum proves that
the symmetric-product map from $X$ to $\mathcal R$ is injective.
By \eqref{eq:rooting}, $\rho\psi$ is precisely that map. Thus $\psi$ is
injective as well.
\end{proof}

For a basis unrooted tree $t$ define the elementary potential
\begin{equation}\label{eq:contraction}
 H_t(V)(q)=\sum_{c:E(t)\to\{1,\ldots,d\}}
               \prod_{a\in\operatorname{Vert}(t)}
                      \partial^{\alpha_a(c)}V(q),
\end{equation}
where $E(t)$ is its edge set and $\alpha_a(c)$ is the multiindex of colors
incident at vertex $a$. For the singleton this is $V$. Extend $H$ linearly
to $\mathcal U$. Differentiation of the finite sum increments the jet at
each possible vertex, giving
\begin{equation}\label{eq:gradient}
 \nabla H_t(V)\cdot\nabla H_v(V)=H_{t\diamond v}(V).
\end{equation}
For potentials $h,k$, our Poisson convention gives
$\{h,T\}=\nabla h\cdot p$ and
$\{k,\{h,T\}\}=\nabla k\cdot\nabla h$.
Induction in \eqref{eq:star} therefore shows
\begin{equation}\label{eq:realization}
 \Phi_{d,V}(x)=H_{\psi(x)}(V)\qquad(x\in X).
\end{equation}

\begin{lemma}\label{lem:jets}
Every nonzero finite linear combination $h\in\mathcal U$ is detected by
$H_h(V)$ for some dimension $d\ge1$ and some polynomial potential $V$.
\end{lemma}
\begin{proof}
Choose a tree $t$ with nonzero coefficient in $h$, and let $r$ be its number
of vertices. As a polynomial in the independent symmetric jet variables
$Z_\alpha=\partial^\alpha V(0)$, its contraction has total jet degree $r$,
so trees with other vertex numbers cannot cancel a degree-$r$ monomial.
For $r=1$, the singleton is the only tree and its monomial is $Z_0$; use $d=1$.

If $r\ge2$, take $d=r-1$ and color its edges with distinct coordinate
indices. Let $\alpha_a$ be the resulting incidence multiindices and set
\[
 M_t=\prod_{a\in\operatorname{Vert}(t)}Z_{\alpha_a}.
\]
Every coordinate index occurs exactly twice. A contraction of another
$r$-vertex tree producing $M_t$ must therefore use each color on exactly one
edge. Its endpoints are the two vertices whose incidence multiindices
contain that color. The multiset of incidence multiindices reconstructs the
colored graph, up to permutations of equal vertices, and hence its unrooted
tree type. In the two-vertex case the two multiindices coincide; exchanging
the endpoints still reconstructs the same tree. Thus only trees isomorphic
to $t$ can contribute $M_t$. Its coefficient in $H_t$ is a positive integer,
since the chosen coloring contributes and all terms have positive coefficients.
Its coefficient in $H_h$ is consequently nonzero.

The nonzero jet polynomial $H_h$ depends on finitely many variables and has
a nonzero real evaluation. One may prove this by induction on the number
of variables: make a nonzero coefficient polynomial nonzero, then avoid
the finitely many roots in the remaining variable.
Arbitrary prescribed jets up to a finite order $D$ are realized by
\[
 V(q)=\sum_{|\alpha|\le D} z_\alpha\frac{q^\alpha}{\alpha!},
 \qquad\partial^\alpha V(0)=z_\alpha.
\]
Taking $D$ at least the largest derivative order needed gives the required
polynomial potential.
\end{proof}

\begin{proof}[Proof of Theorem~\ref{thm:main}]
Let $K$ be the joint kernel over all dimensions and polynomial potentials.
Each realization preserves homogeneous momentum degree. Independence of the
different polynomial degrees in $p$ makes $K$ a graded Lie ideal.
Lemmas~\ref{lem:unrooted} and \ref{lem:jets}, with
\eqref{eq:realization}, show that $K\cap X=0$.
Proposition~\ref{prop:ideal} then gives $K=0$.
Every polynomial potential is smooth, so the joint kernel over smooth
potentials is contained in $K$ and is also zero.
\end{proof}

\appendix
\section{The associative realization}\label{sec:associative}

We recall a short algebraic proof of the fact used in
Lemma~\ref{lem:degreeone}. For an ordinary Lie algebra $L$, choose an ordered
vector-space basis $(e_i)$ and reduce tensor words by
\[
 e_je_i\longmapsto e_ie_j+[e_j,e_i]\qquad(j>i)
\]
in its enveloping algebra. These reductions terminate by tensor length and,
at fixed length, inversion number. Disjoint reductions commute. The only
overlap is a descending three-letter word; the difference of its two
reductions, after reducing shorter terms, is the Jacobi expression and is
zero. Induction on the same decreasing measure resolves every sequence of
reductions. Nondecreasing basis monomials therefore form a basis of the
enveloping algebra, and the length-one terms give an injection $L\to U(L)$.
For the free Lie algebra $L(S)$, the enveloping algebra and the free
associative algebra $\R\langle S\rangle$ have the same universal property:
a map from $S$ to an associative algebra extends uniquely, first as a Lie
map and then as an associative map. They are generator-preservingly
isomorphic. Thus $L(S)$ embeds in $\R\langle S\rangle$.
The argument also applies to infinite $S$, since every word and reduction
is finite.

\section{Exact finite computations}\label{sec:computations}

The accompanying code reconstructs the mechanical-tree matrices and checks
the degree-lowering formulas for every basis element through original-letter
order $18$. A potential tree with $b$ leaves has order $2b-1$, its generator
$y_s$ has order $2b$, and $A$ has order one. Weighted Lyndon words on these
free generators, together with the binary-tree potential basis, provide the
input columns. The marker is a comb with ten leaves. Sparse associative
expansions and tree contractions use integer coefficients; no floating-point
rank threshold is involved.

\begin{center}
\begin{tabular}{@{}rrrrrr@{}}\toprule
Order & 6 & 10 & 14 & 16 & 18\\\midrule
Source dimension & 5 & 39 & 320 & 938 & 2786\\
Mechanical-tree rank & 5 & 39 & 320 & 938 & 2786\\\bottomrule
\end{tabular}
\end{center}
Each rank is computed modulo the prime $1000003$. Full column rank modulo a
prime implies full column rank of the integer matrix over $\mathbb Q$ and
$\R$, since a full-size minor has determinant nonzero modulo that prime.
The stored matrices have $6705$ columns in $91$ order/degree blocks; a
separate elimination replay checks every stored pivot and its integrity hash.
Rebuilding the columns and replaying them are distinct checks.

Across the $6611$ positive-degree columns, the degree-one action is checked
on $362$ elements and the projection/retraction formula
\eqref{eq:euler} on $6249$ elements. All $1019$ associative $y$-words through
the cutoff pass marker decoding. An additional direct sparse-polynomial
implementation in dimension two, with degree-five potentials and seeds
$519$ and $2026$, checks all $40$ columns through order eight for each
potential at $q=0$ and $p=(1,2)$. Exact rational differentiation agrees with
the tree contractions in all $80$ value comparisons. These are pointwise
calibration checks of the implementation. Iterated lowering through order five comprises
$8$ inputs and $17$ steps, all ending in a nonzero potential. These finite
results support implementation correctness; the all-order theorem is proved
in Sections~\ref{sec:lowering}--\ref{sec:potentials}.

The following finite validation procedure takes a cutoff $2\le N\le18$ and a prime
$p_0$; it chooses a marker with $\lfloor N/2\rfloor+1$ leaves and returns
the dimension/rank table and replayable matrices and pivots.
\begin{quote}
\begin{enumerate}
\item Enumerate binary potential trees and weighted Lyndon words of order
$1,\ldots,N$; expand each Lyndon bracket by integer commutators.
\item For every positive-degree column, compute its bracket with the marker.
Check word decoding in degree one and projection followed by retraction in
higher degree.
\item Construct the integer mechanical-tree image of every column, group
columns by order and momentum degree, and eliminate modulo $p_0$ using the
lexicographically first available row as pivot. Reject a dependent column.
Save the integer columns, pivots, labels and integrity hashes for replay.
\end{enumerate}
\end{quote}
If a block
has $C$ columns and $R$ possible rows, dense elimination uses at most
$O(RC^2)$ field operations and $O(RC)$ storage; the implementation uses
sparse dictionaries. These counts concern elimination of an already formed
matrix. Tree generation, word expansion and coefficient bit growth are
additional costs, and no polynomial bound in the order cutoff is asserted.
Caches for sparse intermediate mappings retain at most $4096$ entries per
function; eviction changes neither the arithmetic nor the result.

Reproduction commands, the tested library, the complete order table and
small validation summaries are supplied with the manuscript. The drivers
save full local matrices, lowering traces and source/output hashes in fresh
run directories. The raw matrices can be regenerated from the code.

\begin{thebibliography}{9}
\bibitem{MM}
R.~I. McLachlan and A.~Murua,
\emph{The Lie algebra of classical mechanics},
Journal of Computational Dynamics \textbf{6}(2) (2019), 345--360.
\href{https://doi.org/10.3934/jcd.2019017}{doi:10.3934/jcd.2019017}.
Author preprint: \href{https://arxiv.org/abs/1905.07554}{arXiv:1905.07554}.

\bibitem{BCM}
S.~Blanes, F.~Casas and A.~Murua,
\emph{Splitting methods for differential equations},
Acta Numerica \textbf{33} (2024), 1--161.
\href{https://doi.org/10.1017/S0962492923000077}{doi:10.1017/S0962492923000077}.
Author preprint: \href{https://arxiv.org/abs/2401.01722}{arXiv:2401.01722}.

\bibitem{CL}
F.~Chapoton and M.~Livernet,
\emph{Pre-Lie algebras and the rooted trees operad},
International Mathematics Research Notices \textbf{2001}(8) (2001), 395--408.
\href{https://doi.org/10.1155/S1073792801000198}{doi:10.1155/S1073792801000198}.

\bibitem{BL}
N.~Bergeron and J.-L.~Loday,
\emph{The symmetric operation in a free pre-Lie algebra is magmatic},
Proceedings of the American Mathematical Society \textbf{139}(5) (2011), 1585--1597.
\href{https://doi.org/10.1090/S0002-9939-2010-10813-4}{doi:10.1090/S0002-9939-2010-10813-4}.
Author preprint: \href{https://arxiv.org/abs/1004.1204}{arXiv:1004.1204}.
\end{thebibliography}

\section*{Statement on AI assistance}
OpenAI ChatGPT assisted with literature exploration and the initial proof
development. OpenAI Codex assisted with source verification, proof checking,
manuscript preparation, and implementation and testing of the computational
material. These activities are AI-assisted checks; they do not constitute
independent human peer review. The manuscript is submitted under the author's
name, and independent mathematical review is requested.

\end{document}
